
%(BEGIN_QUESTION)
% Copyright 2006, Tony R. Kuphaldt, released under the Creative Commons Attribution License (v 1.0)
% This means you may do almost anything with this work of mine, so long as you give me proper credit

Hva betyr begrepet {\it rangeability} når det gjelder reguleringsventiler? Hvorfor er denne parameteren viktig for oss?  

\vskip 10pt

Hvilken ventilkarakteristikk har størst rangeability: {\it quick opening}, {\it linear} eller {\it equal percentage}? Begrunn svaret ditt.

\underbar{file i01423}
%(END_QUESTION)





%(BEGIN_ANSWER)

{\it Rangeability} for et måleinstrument viser til den andelen av full skala hvor prosessmålingen kan utføres innenfor oppgitt nøyaktighet. For en reguleringsventil viser ``rangeability'' til den andelen eller det forholdet av full strømning der prosessvæskens strømningsrate kan reguleres pålitelig.
 
Rangeability er viktig fordi den begrenser det praktiske arbeidsområdet for et prosess-strømningssystem. Dersom en ventil har for lav rangeability, kan prosessens strømningsrate ikke reduseres særlig mye fra full strømning og samtidig opprettholde god reguleringskvalitet.

\vskip 10pt

Ventiler med equal percentage-karakteristikk har best rangeability, fordi denne åpningskarakteristikken beholder finst mulig reguleringsgrad mot den nedre delen av ventilslaget sammenlignet med de andre.

Sammenlign en equal percentage-ventil med en quick opening-ventil ved lave posisjonsverdier: quick opening-ventilen vil være svært følsom for spindel-/akselposisjon i nedre område (svært små bevegelser gir store strømningsendringer), mens equal percentage-ventilen beholder gode reguleringsegenskaper mot bunnen av området sitt. Dette gjelder i enda større grad i prosesssystemer med mye karakteristikk-``forvrengning'' (som følge av store endringer i ventilens $\Delta$P over strømningsområdet), der alle iboende ventilkarakteristikker har tendens til å ``bøyes'' mot quick opening-karakteristikk.
 
%(END_ANSWER)





%(BEGIN_NOTES)



%INDEX% Final Control Elements, valve: rangeability

%(END_NOTES)
